\section{Background}\label{sect:back}

\subsection{Conventions}

We use $[m]$ to denote the set $\{1, \dots, m\}$. By $\binom{[m]}{k}$ we mean the set of subsets of $[m]$ of size $k$. For $l \in [m]$, using the notation of \cite{ops}, we use $<_{l}$ to denote the cyclically shifted order on $[m]$ given by \[l <_{l} l+1 <_{l} < \dots <_{l} m-1 <_{l} m <_{l} 1 <_{l} \dots <_{l} l-1 .\] For $r \geqslant 3$, $a_{1} < \dots < a_{r}$ is a \emph{cyclic ordering} if there is an $l \in [n + 2d + 1]$ such that $a_{1} <_{l} \dots <_{l} a_{r}$.

Throughout this paper, we will often need to change the ordering on $[m]$ to $<_{l}$ for some $l$. We will usually do this tacitly, rather than explicitly writing $<_{l}$ in place of $<$. Similarly, we will usually be tacitly using arithmetic modulo $n + 2d + 1$, where this will be the number of vertices of our cyclic polytope. Hence, we shall write things like $j = i - 1$, when we mean $j \equiv i - 1 \pmod{n + 2d + 1}$. Our convention here will also be that our equivalence-class representatives for arithmetic modulo $n + 2d + 1$ will be $[n + 2d + 1]$ rather than $\{0, 1, \dots, n + 2d\}$. This is because we shall label the vertices of our cyclic polytope by $[n + 2d + 1]$.

Finally, given a tuple $A \in [m]^{k + 1}$, we will denote the elements of $A$ by $(a_{0}, a_{1}, \dots, a_{k})$ where $a_{0} <_{l} a_{1} <_{l} \dots <_{l} a_{k}$ in a cyclically shifted order to be determined by the context. By default this will be the usual order on $[m]$. Because we wish sometimes to change this order, we are also tacitly identifying tuples which are the same up to cyclically shifted reordering. The same applies to other letters of the alphabet: the upper case letter denotes the tuple; the lower case letter is used for the entries, which are ordered according to their index, which starts from 0. For example, if we write $A = (1, 3, 5)$, then we have $a_{0} = 1$, $a_{1} = 3$, and $a_{2} = 5$. If we changed the cyclic ordering from $<_{1}$ to $<_{3}$, then we would have $A = (3, 5, 1)$ with $a_{0} = 3$, $a_{1} = 5$, and $a_{2} = 1$.

\subsection{Triangulations of cyclic polytopes}\label{sect:back:cyc}

Cyclic polytopes should be thought of as higher-dimensional analogues of convex polygons. General introductions to this class of polytopes can be found in \cite[Lecture 0]{ziegler}, \cite[4.7]{gruenbaum}, \cite[Section 6.1]{lrs} and \cite[Chapter VI]{bar}. Let $n \geqslant 0$ and $\delta \geqslant 1$ and consider the curve defined by $p_{t}=(t, t^{2}, \dots , t^{\delta}) \subset \mathbb{R}^{\delta}$ for $t \in \mathbb{R}$, which is known as the \emph{moment curve}. Choose $t_{1}, t_{2}, \dots , t_{n + \delta + 1} \in \mathbb{R}$ such that $t_{1} < t_{2} < \dots < t_{n + \delta + 1}$. The convex hull $\conv\{ p_{t_{1}}, p_{t_{2}}, \dots , p_{t_{n + \delta + 1}} \}$ is a \emph{cyclic polytope} $C(n + \delta + 1, \delta)$. A (\emph{geometric}) \emph{triangulation} of a cyclic polytope $C(n + \delta + 1, \delta)$ is a collection of (geometric) $\delta$-simplices whose interiors are pairwise disjoint and whose union is $C(n + \delta + 1, \delta)$.

A \emph{facet} of $C(n + \delta + 1, \delta)$ is a face of codimension one. A \emph{circuit} of a cyclic polytope $C(n + \delta + 1, \delta)$ is a pair, $(Z_{+}, Z_{-})$, of sets of vertices of $C(n + \delta + 1,\delta)$ which are inclusion-minimal with respect to the property $\mathrm{conv}(Z_{+}) \cap \mathrm{conv}(Z_{-}) \neq \emptyset$, in the sense that if $Z'_{+} \subseteq Z_{+}$ and $Z'_{-} \subseteq Z_{-}$ with $\mathrm{conv}(Z'_{+}) \cap \mathrm{conv}(Z'_{-}) \neq \emptyset$, then $Z'_{+} = Z_{+}$ and $Z'_{-} = Z_{-}$.

The circuits and facets of a cyclic polytope are independent of its particular geometric realisation given by the choice of points on the moment curve, by \cite{gale,breen}. This implies that whether or not a collection of ordered $(\delta + 1)$-tuples of $[n + \delta + 1]$ forms a triangulation of $C(n + \delta + 1, \delta)$ is likewise independent of the particular geometric realisation. Hence, in this paper we primarily consider triangulations combinatorially: a \emph{combinatorial triangulation} of $C(n + \delta + 1, \delta)$ is a collection of ordered $(\delta + 1)$-tuples with entries in $[n + \delta + 1]$ which gives a geometric triangulation of $C(n + \delta + 1, \delta)$. Unless otherwise specified, when we write `triangulation' we mean a combinatorial triangulation. Likewise, if we say that $A$ is a $k$-simplex, we mean that $A$ is an ordered $(k + 1)$-tuple with entries in $[n + \delta + 1]$. We write $|A|$ to denote the geometric realisation of $A$ as a geometric simplex which is the convex hull of $k + 1$ points on the moment~curve. 

In this paper we are exclusively interested in triangulations of even-dimensional cyclic polytopes. These were given an elegant combinatorial description in \cite{ot}, which we now explain. This combinatorial description was in turn used to connect triangulations of even-dimensional cyclic polytopes to representation theory of algebras. We will sometimes comment on these connections, but this paper does not require the reader to be familiar with representation theory.

A $d$-simplex of a triangulation of $C(n + 2d + 1, 2d)$ is \emph{internal} if it does not lie within a facet of $C(n + 2d + 1, 2d)$. It is clear that a triangulation of a convex polygon is determined by the arcs of the triangulation; similarly, a triangulation of $C(n + 2d + 1,2d)$ is determined by the internal $d$-simplices of the triangulation, by a theorem of Dey \cite{dey}. For brevity, we refer to an internal $d$-simplex of a triangulation $\mathcal{T}$ of $C(n + 2d + 1, 2d)$ as a \emph{$d$-arc} of $\mathcal{T}$. We write $\darcs{\mathcal{T}}$ for the set of $d$-arcs of $\mathcal{T}$. We say that a $(d + 1)$-simplex of $\mathcal{T}$ is \emph{interior} if all of its facets are $d$-arcs.

In order to recall the description of the triangulations of $C(n + 2d + 1,2d)$ from \cite{ot}, and for use later in the paper, we denote the following sets.
\begin{align*}
\derset{n + 2d + 1}{d} &= \set{(a_{0},\dots,a_{d}) \in \mathbb{Z}^{d+1} \st \parbox{5.5cm}{\begin{center}$\forall i \in \{0, 1, \dots ,d - 1\}, a_{i + 1} \geqslant a_{i} + 2 $\\  and  $a_{d} + 2 \leqslant a_{0} + n + 2d + 1$\end{center}}}; \\
\nonconsec{n + 2d + 1}{d} &= \derset{n + 2d + 1}{d} \cap [n + 2d + 1]^{d + 1}.
\end{align*}

\begin{rema}
Algebraically, $\derset{n + 2d + 1}{d}$ labels the $d$-cluster-tilting subcategory of the derived category of $A_{n}^{d}$, whilst $\nonconsec{n + 2d + 1}{d}$ labels its cluster category \cite{ot}. Here $A_{n}^{d}$ is the higher Auslander algebra of type~$A$; see \cite{iy-clus}.
\end{rema}

A $d$-simplex $A \in [n + 2d + 1]^{d + 1}$ is a $d$-arc in $C(n + 2d + 1, 2d)$ if and only if $A \in \nonconsec{n + 2d + 1}{d}$. A $d$-simplex $A$ and a $d$-simplex $B$ are \emph{intertwining} if \[ a_{0} < b_{0} < a_{1} < b_{1} < \dots < a_{d} < b_{d}\] is a cyclic ordering, in which case we write $A \swr B$. Note that $A$ and $B$ being intertwining is therefore independent of the cyclically shifted order $<_{l}$ we may choose on $[n + 2d + 1]$. A collection of $d$-simplices is called \emph{non-intertwining} if no pair of its elements are intertwining. The circuits of the cyclic polytope $C(n + 2d + 1, 2d)$ are pairs of intertwining $d$-arcs.

\begin{theorem}[{\cite[Theorem 2.3, Theorem 2.4]{ot}}]
There is a bijection between elements of $\nonconsec{n + 2d + 1}{d}$ and $d$-arcs of $C(n + 2d + 1, 2d)$, which induces a bijection between non-intertwining collections of $\binom{n + d - 1}{d}$ $d$-simplices in $\nonconsec{n + 2d + 1}{d}$ and triangulations of $C(n + 2d + 1, 2d)$, via sending \[\mathcal{T} \mapsto \darcs{\mathcal{T}}.\]
\end{theorem}

Triangulations of cyclic polytopes can be \emph{mutated} by operations known as \emph{bistellar flips} \cite[Section~2.4]{lrs}. The following theorem can be taken as a definition of this operation in the case of triangulations of $C(n + 2d + 1, 2d)$.

\begin{theorem}[{\cite[Theorem~4.1]{ot}}]\label{thm:flip}
Triangulations $\mathcal{T}$ and $\mathcal{T}'$ of $C(n + 2d + 1,2d)$ are bistellar flips of each other if and only if $\darcs{\mathcal{T}}$ and $\darcs{\mathcal{T}'}$ have all but one element in common.
\end{theorem}
